\subsubsection{AN TOÀN KHI SỬ DỤNG THIẾT BỊ THÍ NGHIỆM}
\immini{\paragraph{Sử dụng các thiết bị điện}
\par Quan sát kĩ các kí hiệu và nhãn thông số trên thiết bị để sử dụng đúng chức năng, đúng yêu cầu kĩ thuật.
\paragraph{Sử dụng các thiết bị nhiệt và thuỷ tinh}
\par Các thiết bị đun nóng có thể gây bỏng với người sử dụng, gây nứt, vỡ các bộ phận làm bằng thủy tinh.
}{\includegraphics[width=4cm]{images/adapter}}
\paragraph{Sử dụng các thiết bị quang học}
\par Các thiết bị quang học rất dễ mốc, xước, nứt, vỡ và dính bụi bẩn, làm ảnh hưởng đến đường truyền tia sáng và sai lệch kết quả thí nghiệm.
\begin{center}
\includegraphics[width=1\linewidth]{images/LT-2-1S}
\end{center}
\subsubsection{NGUY CƠ MẤT AN TOÀN TRONG SỬ DỤNG THIẾT BỊ THÍ NGHIỆM VẬT LÍ}
\paragraph{Nguy cơ gây nguy hiểm cho người sử dụng}
\par Khi tiến hành thí nghiệm cần tuân thủ nghiêm ngặt các quy định trong phòng thực hành và hướng dẫn của giáo viên.
\par Một số thao tác có thể gây mất an toàn trong phòng thí nghiệm:
% TODO: \usepackage{graphicx} required
\begin{center}
\includegraphics[width=1\linewidth]{images/LT-2-2S}
\end{center}
\immini{
\paragraph{Nguy cơ hỏng thiết bị đo điện}
\par Khi sử dụng các thiết bị đo điện cần chọn đúng thang đo, không nhầm lẫn khi thao tác để đảm bảo an toàn cho thiết bị đo.
\par Khi sử dụng đồng hồ đo điện đa năng (thiết bị đo điện với các chức năng chính là đo điện trở, đo hiệu điện thế và đo dòng điện AC, DC), cần lưu ý:
\begin{itemize}
\item Chọn chức năng và thang đo phù hợp.
\item Cắm dây đo vào chốt cắm phù hợp với chức năng đo.
\end{itemize}
\paragraph{Nguy cơ cháy nổ trong phòng thực hành}
Khi tiến hành thí nghiệm với các thiết bị điện và những hoá chất, chất dễ cháy nổ trong phòng thực hành cần tuân thủ quy tắc an toàn, nhất là những quy tắc an toàn về phòng cháy chữa cháy và an toàn khi sử dụng hoá chất dễ cháy, nổ}{% TODO: \usepackage{graphicx} required
\includegraphics[width=5cm]{images/LT-2-3}
}

\subsubsection{QUY TẮC AN TOÀN TRONG PHÒNG THỰC HÀNH}
\begin{itemize}
\item Đọc kĩ hướng dẫn sử dụng thiết bị và quan sát các chỉ dẫn, các kí hiệu trên các thiết bị thí nghiệm.
\item Kiểm tra cẩn thận thiết bị, phương tiện, dụng cụ thí nghiệm trước khi sử dụng.
\item Chỉ tiến hành thí nghiệm khi được sự cho phép của giáo viên hướng dẫn thí nghiệm.
\item Tắt công tắc nguồn thiết bị điện trước khi cắm hoặc tháo thiết bị điện.
\item Chỉ cắm phích / giắc cắm của thiết bị điện vào ổ cắm khi hiệu điện thế của nguồn điện tương ứng với hiệu điện thế định mức của dụng cụ.
\item Phải bố trí dây điện gọn gàng, không bị vướng khi qua lại.
\item Không tiếp xúc trực tiếp với các vật và các thiết bị thí nghiệm có nhiệt độ cao khi không có dụng cụ bảo hộ.
\item Không để nước cũng như các dung dịch dẫn điện, dung dịch dễ cháy gần thiết bị điện.
\item Giữ khoảng cách an toàn khi tiến hành thí nghiệm nung nóng các vật, thí nghiệm có các vật bắn ra, tia laser.
\item Phải vệ sinh, sắp xếp gọn gàng các thiết bị và dụng cụ thí nghiệm, bỏ chất thải thí nghiệm vào đúng nơi quy định sau khi tiến hành thí nghiệm.
\end{itemize}
\begin{center}
\includegraphics[width=0.2\linewidth]{"images/2 (2)"}
\includegraphics[width=0.2\linewidth]{"images/2 (4)"}
\includegraphics[width=0.2\linewidth]{"images/2 (5)"}
\includegraphics[width=0.2\linewidth]{"images/2 (3)"}
\end{center}
\begin{center}
\includegraphics[width=0.2\linewidth]{"images/2 (6)"}
\includegraphics[width=0.2\linewidth]{"images/2 (8)"}
\includegraphics[width=0.21\linewidth]{"images/2 (1)"}
\includegraphics[width=0.2\linewidth]{"images/2 (7)"}
\end{center}